\title{Επισκόπηση Λειτουργίας Συστήματος Μεταφοράς Ηλεκτρικής Ενέργειας}

\section{Θεωρία}
\label{kef-01-theoria}

Το κεφάλαιο αυτό απαντά σε ένα ερώτημα που φαίνεται απλό αλλά δεν είναι: \emph{πώς
κρατάει κανείς ένα σύστημα ηλεκτρικής ενέργειας σε λειτουργία}. Θα δούμε ότι η απάντηση
είναι ένας κλειστός βρόχος τεσσάρων σταδίων — μετράμε το σύστημα, μετατρέπουμε τις
μετρήσεις σε αξιόπιστη εικόνα, αξιολογούμε την ασφάλεια αυτής της εικόνας, και δρούμε.
Οι ενότητες που ακολουθούν εξετάζουν με τη σειρά αυτά τα στάδια. Το ζητούμενο δεν είναι
να απομνημονευτούν τα ακρωνύμια, αλλά να γίνει κατανοητό γιατί κάθε στάδιο υπάρχει και
τι θα πήγαινε στραβά χωρίς αυτό.

\subsection{Ο ρόλος του Συστήματος Μεταφοράς}

Το Σύστημα Μεταφοράς Ηλεκτρικής Ενέργειας (ΣΜΗΕ) είναι το δίκτυο υπερυψηλής και υψηλής
τάσης που διασυνδέει τους σταθμούς παραγωγής με τα κέντρα κατανάλωσης και με τα γειτονικά
συστήματα. Στην Ελλάδα λειτουργεί στα επίπεδα των 400 kV, 150 kV και 66 kV και τη
διαχείρισή του ασκεί ο Ανεξάρτητος Διαχειριστής Μεταφοράς Ηλεκτρικής Ενέργειας (ΑΔΜΗΕ),
ο οποίος είναι μέλος του ENTSO-E και τμήμα της σύγχρονης περιοχής της Ηπειρωτικής Ευρώπης
(Continental Europe).

Γιατί μεταφέρουμε ενέργεια σε τόσο υψηλή τάση; Για δεδομένη μεταφερόμενη ισχύ
$P = \sqrt{3}\,V I \cos\varphi$, το ρεύμα είναι αντιστρόφως ανάλογο της τάσης, ενώ οι
απώλειες στους αγωγούς είναι ανάλογες του $I^2$. Διπλασιασμός της τάσης υποτετραπλασιάζει
επομένως τις απώλειες μεταφοράς. Αυτός είναι και ο λόγος που το σύστημα οργανώνεται
ιεραρχικά: υπερυψηλή τάση για τις μεγάλες αποστάσεις, διαδοχικοί υποβιβασμοί καθώς
πλησιάζουμε τον καταναλωτή.

\subsubsection*{Οι τρεις ταυτόχρονες απαιτήσεις}

Ο Διαχειριστής Συστήματος Μεταφοράς (Transmission System Operator, TSO) καλείται να
ικανοποιεί ταυτόχρονα τρεις απαιτήσεις. Η ένταση μεταξύ τους είναι που κάνει το πρόβλημα
ενδιαφέρον:

\begin{itemize}
\item \textbf{Ισοζύγιο ισχύος.} Η ηλεκτρική ενέργεια δεν αποθηκεύεται σε σημαντικές
      ποσότητες μέσα στο ίδιο το σύστημα. Παραγωγή και ζήτηση πρέπει να εξισώνονται
      \emph{συνεχώς}. Όταν δεν εξισώνονται, η διαφορά καλύπτεται προσωρινά από την
      κινητική ενέργεια των στρεφόμενων μαζών, και το τίμημα είναι μεταβολή της
      συχνότητας — ένα σύστημα που καταναλώνει περισσότερο από όσο παράγει επιβραδύνεται
      κυριολεκτικά.
\item \textbf{Ασφάλεια λειτουργίας.} Το σύστημα πρέπει να παραμένει εντός ορίων (θερμικά
      όρια στοιχείων, όρια τάσεων, όρια ευστάθειας) όχι μόνο τώρα, αλλά και μετά από κάθε
      πιθανή μεμονωμένη διαταραχή. Αυτή η διάκριση μεταξύ «τώρα» και «μετά από» είναι
      θεμελιώδης και οδηγεί στο κριτήριο N-1 που θα δούμε αμέσως παρακάτω.
\item \textbf{Ποιότητα.} Η συχνότητα και οι τάσεις πρέπει να διατηρούνται εντός
      προδιαγεγραμμένων περιθωρίων, τα οποία για την Ηπειρωτική Ευρώπη ορίζονται
      ρυθμιστικά και όχι κατά την κρίση του Διαχειριστή.
\end{itemize}

Η δυσκολία είναι ότι οι απαιτήσεις αυτές ζουν σε \emph{διαφορετικές χρονικές κλίμακες},
που εκτείνονται σε δώδεκα τάξεις μεγέθους. Ο Πίνακας~\ref{tab:chronikes-klimakes}
συνοψίζει την εικόνα και λειτουργεί ως χάρτης για το υπόλοιπο κεφάλαιο.

\begin{table}[!htbp]
\centering
\begin{tabular}{lll}
\hline
Κλίμακα & Φαινόμενο / ενέργεια & Μηχανισμός \\
\hline
μs -- ms   & κύματα, σφάλματα           & προστασία, διακόπτες \\
ms -- s    & ηλεκτρομηχανικές ταλαντώσεις & αδράνεια, PSS, FACTS \\
s -- 30 s  & πτώση συχνότητας            & πρωτεύουσα ρύθμιση (FCR) \\
30 s -- 15 min & αποκατάσταση συχνότητας & δευτερεύουσα/τριτεύουσα (aFRR, mFRR) \\
min -- h   & αναπρογραμματισμός          & αγορά εξισορρόπησης, ενδοημερήσια \\
h -- ημέρα & κατανομή φορτίου            & αγορά επόμενης ημέρας \\
εβδομάδες -- έτη & επάρκεια, συντηρήσεις & προγραμματισμός, επενδύσεις \\
\hline
\end{tabular}
\caption{Χρονικές κλίμακες λειτουργίας του ΣΜΗΕ. Κάθε γραμμή αντιστοιχεί σε διαφορετικό
μηχανισμό ελέγχου· το κεφάλαιο αυτό εστιάζει από τα δευτερόλεπτα και πάνω.}
\label{tab:chronikes-klimakes}
\end{table}

\subsection{Καταστάσεις λειτουργίας και το κριτήριο N-1}

Ένα σύστημα δεν είναι απλώς «εντάξει» ή «χαλασμένο». Η καθιερωμένη σήμερα εικόνα οφείλεται
στον Dy Liacco, ο οποίος στα τέλη της δεκαετίας του 1960 διατύπωσε την ιδέα ότι ο έλεγχος
οφείλει να \emph{προσαρμόζεται} στην εκάστοτε κατάσταση ασφαλείας. Το ευρωπαϊκό κανονιστικό
πλαίσιο (System Operation Guideline, Κανονισμός ΕΕ 2017/1485, Άρθρο 18) κωδικοποιεί την
ιδέα αυτή σε πέντε καταστάσεις:

\begin{enumerate}
\item \textbf{Κανονική} (normal). Όλα τα μεγέθη εντός ορίων \emph{και} το κριτήριο N-1
      ικανοποιείται. Ο έλεγχος είναι \emph{προληπτικός}: ο Διαχειριστής προετοιμάζεται για
      διαταραχές που δεν έχουν συμβεί.
\item \textbf{Επιφυλακής} (alert). Τα μεγέθη παραμένουν εντός ορίων, αλλά έχει εξαντληθεί
      το περιθώριο εφεδρείας ή παραβιάζεται το N-1. Το σύστημα λειτουργεί, όμως μια επόμενη
      διαταραχή θα το βγάλει εκτός ορίων.
\item \textbf{Έκτακτης ανάγκης} (emergency). Παραβιάζεται τουλάχιστον ένα όριο ασφαλούς
      λειτουργίας. Ο έλεγχος γίνεται \emph{διορθωτικός} και επείγων.
\item \textbf{Συσκότισης} (blackout). Εκτεταμένη απώλεια φορτίου ή κατάρρευση τμήματος
      του συστήματος.
\item \textbf{Αποκατάστασης} (restoration). Ενεργοποιείται το σχέδιο επανατροφοδότησης.
\end{enumerate}

Η μετάβαση από την κανονική στην κατάσταση επιφυλακής \emph{δεν} συνοδεύεται από κάποια
ορατή αλλαγή στα μετρούμενα μεγέθη — τίποτα δεν υπερβαίνει κάποιο όριο. Είναι μια αλλαγή
που υπάρχει μόνο στους υπολογισμούς του Διαχειριστή. Γι' αυτό ακριβώς χρειάζεται μια
αξιόπιστη εικόνα του συστήματος: χωρίς αυτήν, η διαφορά μεταξύ ασφαλούς και επικίνδυνης
λειτουργίας είναι αόρατη.

\subsubsection*{Το κριτήριο N-1}

Το \textbf{κριτήριο N-1} είναι ο πρακτικός ορισμός της ασφάλειας: από τα $N$ στοιχεία του
συστήματος, η απώλεια οποιουδήποτε \emph{ενός} — γραμμής, μετασχηματιστή, μονάδας παραγωγής,
ζυγού — πρέπει να μπορεί να απορροφηθεί χωρίς παραβίαση ορίων και χωρίς αλυσιδωτές
αποσυνδέσεις.

Ο έλεγχός του γίνεται με \emph{ανάλυση ενδεχομένων} (contingency analysis): για κάθε
πιθανό ενδεχόμενο εκτελείται μια ροή ισχύος με το αντίστοιχο στοιχείο εκτός, και ελέγχεται
αν κάποιο μέγεθος υπερβαίνει τα όριά του. Σε ένα σύστημα με μερικές εκατοντάδες κλάδους,
αυτό σημαίνει εκατοντάδες ροές ισχύος που πρέπει να ολοκληρωθούν μέσα σε λίγα λεπτά.

Δύο παρατηρήσεις που αξίζει να συγκρατηθούν:

\begin{itemize}
\item Το N-1 είναι \emph{ντετερμινιστικό} κριτήριο: δεν σταθμίζει την πιθανότητα κάθε
      ενδεχομένου. Ένα σπάνιο και ένα συχνό σφάλμα αντιμετωπίζονται εξίσου. Είναι απλό και
      ελεγχόμενο, αλλά συντηρητικό σε κάποιες περιπτώσεις και ανεπαρκές σε άλλες — εξ ου
      και το ενδιαφέρον για πιθανοτικά κριτήρια ασφαλείας.
\item Η ανάλυση ενδεχομένων απαιτεί ως \emph{είσοδο} την τρέχουσα κατάσταση του συστήματος.
      Αν αυτή η είσοδος είναι λάθος, όλα τα συμπεράσματα περί ασφάλειας είναι λάθος. Αυτό
      καθιστά την εκτίμηση κατάστασης (§\ref{sec:se}) όχι απλώς χρήσιμη, αλλά προαπαιτούμενη.
\end{itemize}

\subsection{Εποπτεία και έλεγχος: SCADA και EMS}

\subsubsection*{Η αλυσίδα μέτρησης}

Η βασική υποδομή τηλεμέτρησης και τηλεχειρισμού του ΣΜΗΕ είναι το σύστημα \textbf{SCADA}
(Supervisory Control and Data Acquisition). Η αλυσίδα από το φυσικό μέγεθος έως την οθόνη
του χειριστή έχει τέσσερις κρίκους:

\begin{enumerate}
\item \textbf{Μετασχηματιστές οργάνων.} Οι μετασχηματιστές τάσης (VT) και έντασης (CT)
      υποβιβάζουν τα μεγέθη του δικτύου σε επίπεδα που χειρίζονται τα ηλεκτρονικά. Εισάγουν
      το πρώτο σφάλμα, τυπικά κλάσης 0,2 έως 1\%.
\item \textbf{Μεταλλάκτες και IED.} Υπολογίζουν τα παράγωγα μεγέθη (ενεργό και άεργο ισχύ,
      ενεργό τιμή τάσης) από τα στιγμιαία δείγματα.
\item \textbf{RTU.} Οι Απομακρυσμένες Τερματικές Μονάδες (Remote Terminal Units)
      συγκεντρώνουν τα δεδομένα του υποσταθμού και τα αποστέλλουν στο Κέντρο Ελέγχου.
\item \textbf{Τηλεπικοινωνιακό δίκτυο.} Μεταφέρει τα δεδομένα με τυποποιημένα πρωτόκολλα,
      με πλέον διαδεδομένα τα IEC 60870-5-101/104 και το DNP3, ενώ εντός του υποσταθμού
      κυριαρχεί πλέον το IEC 61850.
\end{enumerate}

Συλλέγονται δύο κατηγορίες δεδομένων:

\begin{itemize}
\item \textbf{Αναλογικές μετρήσεις}: ενεργός και άεργος ροή ισχύος στα άκρα των γραμμών και
      των μετασχηματιστών, εγχεόμενη ισχύς σε ζυγούς, μέτρα τάσεων, ρεύματα κλάδων.
\item \textbf{Ψηφιακές ενδείξεις}: θέσεις διακοπτών και αποζευκτών, θέσεις λήψεων
      μετασχηματιστών, σήματα συναγερμών και λειτουργίας προστασιών. Αυτές καθορίζουν την
      \emph{τοπολογία} — ποια στοιχεία είναι εν λειτουργία και πώς συνδέονται μεταξύ τους.
\end{itemize}

Ο ρυθμός σάρωσης ενός συμβατικού SCADA είναι τυπικά της τάξης των λίγων δευτερολέπτων
(συνήθως 2--10 s).

\subsubsection*{Δύο περιορισμοί που καθορίζουν τα πάντα}

Δύο χαρακτηριστικά του SCADA είναι κρίσιμα για όσα ακολουθούν, και αξίζει να τα κρατήσει
κανείς:

\begin{itemize}
\item Οι μετρήσεις \textbf{δεν είναι χρονικά συγχρονισμένες} μεταξύ τους. Συλλέγονται σε
      διαφορετικές στιγμές εντός του κύκλου σάρωσης, οπότε το σύνολο των μετρήσεων αποτελεί
      μια \emph{ημι-στατική} απεικόνιση (snapshot) και όχι στιγμιότυπο. Ένα σύνολο
      μετρήσεων που έχει «απλωθεί» σε πέντε δευτερόλεπτα δεν αντιστοιχεί σε καμία
      πραγματική στιγμή του συστήματος. Αυτό δεν είναι πρόβλημα όσο το σύστημα μεταβάλλεται
      αργά — και είναι σοβαρό πρόβλημα κατά τη διάρκεια μιας διαταραχής.
\item Οι μετρήσεις είναι \textbf{θορυβώδεις και ενίοτε εσφαλμένες}, ενώ ορισμένα μεγέθη
      δεν μετρώνται καθόλου. Χαρακτηριστικά, οι \emph{γωνίες} των τάσεων — που είναι το ήμισυ
      της κατάστασης του συστήματος — απουσιάζουν εντελώς από το συμβατικό SCADA.
\end{itemize}

Ένας εσφαλμένος μετασχηματιστής έντασης, μια RTU με λάθος πολικότητα, ένας διακόπτης του
οποίου η ένδειξη θέσης έχει κολλήσει: όλα αυτά παράγουν δεδομένα που μοιάζουν απολύτως
φυσιολογικά. Το σύστημα εποπτείας δεν έχει τρόπο να τα ξεχωρίσει \emph{ανά μέτρηση}. Θα
δούμε ότι μπορεί να τα ξεχωρίσει μόνο συλλογικά, εκμεταλλευόμενο τον πλεονασμό.

\subsubsection*{Το EMS και η σειρά των λειτουργιών}

Το \textbf{EMS} (Energy Management System) είναι το λογισμικό που μετατρέπει τα ακατέργαστα
δεδομένα σε αποφάσεις. Η σειρά με την οποία εκτελούνται οι λειτουργίες του δεν είναι
αυθαίρετη — κάθε βήμα τροφοδοτεί το επόμενο:

\begin{enumerate}
\item \textbf{Επεξεργαστής τοπολογίας} (topology processor). Από τις θέσεις των διακοπτών
      συνθέτει το ηλεκτρικό μοντέλο ζυγών-κλάδων (bus-branch model). Απαντά στο ερώτημα
      «πώς είναι συνδεδεμένο το δίκτυο αυτή τη στιγμή;».
\item \textbf{Εκτίμηση κατάστασης}. Απαντά στο «ποια είναι η ηλεκτρική κατάσταση αυτού του
      δικτύου;».
\item \textbf{Ανάλυση ενδεχομένων}. Απαντά στο «είναι ασφαλής αυτή η κατάσταση;».
\item \textbf{Βέλτιστη ροή ισχύος} (Optimal Power Flow). Απαντά στο «αν δεν είναι, τι
      αλλάζουμε και με ποιο κόστος;».
\item \textbf{Αυτόματος έλεγχος παραγωγής} (Automatic Generation Control). Υλοποιεί συνεχώς
      τη ρύθμιση συχνότητας και τις διασυνοριακές ανταλλαγές.
\end{enumerate}

Ένα λάθος στο βήμα 1 ή 2 διαδίδεται σιωπηλά σε όλα τα υπόλοιπα.

\subsection{Συγχρονισμένες μετρήσεις φασιθετών (PMU) και συστήματα WAMS}

\subsubsection*{Γιατί έχει σημασία η γωνία}

Θυμηθείτε τη σχέση μεταφοράς ενεργού ισχύος μεταξύ δύο ζυγών που συνδέονται με
επαγωγική γραμμή αντίδρασης $X$:

\begin{equation}
P_{12} = \frac{V_1 V_2}{X}\,\sin(\delta_1 - \delta_2)
\label{eq:powerflow}
\end{equation}

Η ροή ισχύος καθορίζεται από τη \emph{διαφορά γωνιών}. Η γωνία δεν είναι λοιπόν ένα
βοηθητικό μαθηματικό μέγεθος: είναι η κινητήρια δύναμη της ροής ισχύος και ο πλέον άμεσος
δείκτης καταπόνησης του δικτύου. Το ότι το συμβατικό SCADA δεν τη μετρά είναι σοβαρή
έλλειψη.

\subsubsection*{Η αρχή λειτουργίας}

Οι \textbf{Μονάδες Μέτρησης Φασιθετών} (Phasor Measurement Units, PMU) αίρουν αυτόν τον
περιορισμό. Μια PMU υπολογίζει τον φασιθέτη — μέτρο \emph{και} γωνία — της τάσης και του
ρεύματος, και τον συνοδεύει με χρονοσήμανση από κοινή χρονική αναφορά UTC, τυπικά μέσω GPS.

Το κρίσιμο σημείο είναι το εξής: η γωνία ενός φασιθέτη έχει νόημα μόνο ως προς κάποια
αναφορά. Επειδή όλες οι PMU αναφέρονται στην ίδια χρονική βάση, οι γωνίες που μετρούν είναι
άμεσα συγκρίσιμες σε ολόκληρο το σύστημα, ακόμη και σε αποστάσεις εκατοντάδων χιλιομέτρων.
Εξ ου και ο όρος \emph{συγχρονοφασιθέτες} (synchrophasors). Για να έχει νόημα η σύγκριση σε
σύστημα 50 Hz χρειάζεται ακρίβεια χρονισμού καλύτερη του μικροδευτερολέπτου — σφάλμα 1 μs
αντιστοιχεί σε σφάλμα γωνίας $360^\circ \times 50 \times 10^{-6} = 0{,}018^\circ$.

\subsubsection*{Προδιαγραφές}

Τα βασικά τεχνικά χαρακτηριστικά καθορίζονται από το πρότυπο IEC/IEEE 60255-118-1:2018, το
οποίο διαδέχθηκε τη σειρά IEEE C37.118:

\begin{itemize}
\item \textbf{Ρυθμός αναφοράς}: τυπικά 10, 25 ή 50 πλαίσια ανά δευτερόλεπτο για σύστημα
      50 Hz, με δυνατότητα υψηλότερων ρυθμών.
\item \textbf{Ακρίβεια}: το συνολικό σφάλμα φασιθέτη (Total Vector Error, TVE) δεν πρέπει
      να υπερβαίνει το 1\% στη μόνιμη κατάσταση. Ορίζεται ως το μέτρο της διαφοράς του
      μετρούμενου από τον πραγματικό φασιθέτη, κανονικοποιημένο ως προς τον πραγματικό, και
      επομένως ενοποιεί σε ένα μέγεθος τα σφάλματα μέτρου και γωνίας. Ορίζονται επιπλέον
      το σφάλμα συχνότητας (FE) και το σφάλμα ρυθμού μεταβολής συχνότητας (RFE).
\item \textbf{Κλάσεις}: η κλάση P (protection) βελτιστοποιείται για ταχεία απόκριση, με
      μικρή καθυστέρηση φίλτρου, και η κλάση M (measurement) για ακρίβεια υπό παρουσία
      αρμονικών και παρεμβολών. Η επιλογή είναι συμβιβασμός: δεν υπάρχει φίλτρο που να
      είναι ταυτόχρονα ταχύ και ισχυρά απορριπτικό.
\end{itemize}

\begin{table}[!htbp]
\centering
\begin{tabular}{lll}
\hline
 & SCADA & PMU / WAMS \\
\hline
Ρυθμός           & 2--10 s ανά σάρωση & 10--50 πλαίσια/s \\
Συγχρονισμός     & όχι                & ναι (UTC, GPS) \\
Γωνία τάσης      & δεν μετράται       & μετράται άμεσα \\
Φύση μετρήσεων   & μη γραμμικές ως προς την κατάσταση & γραμμικές \\
Κάλυψη           & εκτεταμένη         & επιλεκτική (κόστος) \\
Τυπική χρήση     & εκτίμηση κατάστασης, τηλεχειρισμός & δυναμικά φαινόμενα \\
\hline
\end{tabular}
\caption{Σύγκριση συμβατικής και συγχρονισμένης μέτρησης. Οι δύο τεχνολογίες είναι
συμπληρωματικές και όχι ανταγωνιστικές.}
\label{tab:scada-pmu}
\end{table}

\subsubsection*{Από τη μέτρηση στην εφαρμογή}

Οι μετρήσεις συγκεντρώνονται ιεραρχικά σε Συμπυκνωτές Δεδομένων Φασιθετών (Phasor Data
Concentrators, PDC), οι οποίοι ευθυγραμμίζουν χρονικά τα πλαίσια από διαφορετικές PMU και
συνθέτουν ένα \textbf{Σύστημα Εποπτείας Ευρείας Περιοχής} (Wide Area Monitoring System,
WAMS). Οι χαρακτηριστικές εφαρμογές είναι:

\begin{itemize}
\item παρακολούθηση ηλεκτρομηχανικών ταλαντώσεων μεταξύ περιοχών, με συχνότητες τυπικά
      0,1--1 Hz, οι οποίες είναι πρακτικά αόρατες σε ρυθμό σάρωσης δευτερολέπτων·
\item μέτρηση του ρυθμού μεταβολής της συχνότητας και εκτίμηση της διαθέσιμης αδράνειας·
\item ανίχνευση νησιδοποίησης και επιτήρηση γωνιακής απόκλισης μεταξύ περιοχών·
\item επικύρωση μοντέλων μετά από διαταραχές, συγκρίνοντας την καταγεγραμμένη με την
      προσομοιωμένη απόκριση·
\item ενίσχυση — ή, σε πλήρη κάλυψη, αντικατάσταση — της κλασικής εκτίμησης κατάστασης.
\end{itemize}

\subsection{Διατάξεις ελεγχόμενης μεταφοράς: FACTS}

Ενώ SCADA και PMU \emph{παρατηρούν} το σύστημα, οι διατάξεις \textbf{FACTS} (Flexible AC
Transmission Systems) το \emph{ελέγχουν} ενεργά.

Η \eqref{eq:powerflow} δείχνει ακριβώς τι μπορεί να ελεγχθεί. Η ροή ισχύος εξαρτάται από
τρία και μόνο μεγέθη: τις τάσεις $V_1, V_2$, την αντίδραση $X$ και τη διαφορά γωνιών
$\delta_1 - \delta_2$. Σε ένα κλασικό δίκτυο εναλλασσομένου ρεύματος και τα τρία είναι
ουσιαστικά δεδομένα — η ισχύς ρέει όπου την οδηγεί η φυσική, όχι όπου θα θέλαμε. Οι
διατάξεις FACTS είναι μετατροπείς ηλεκτρονικών ισχύος που καθιστούν αυτά τα μεγέθη
ελεγχόμενα, ταχύτατα και συνεχώς. Κατατάσσονται ανάλογα με το ποιο από αυτά επηρεάζουν:

\begin{itemize}
\item \textbf{Παράλληλης σύνδεσης} (ελέγχουν το $V$): ο Στατός Αντισταθμιστής Αέργου Ισχύος
      (Static Var Compensator, SVC) και ο STATCOM. Ελέγχουν την τάση του ζυγού μέσω έγχυσης
      ή απορρόφησης αέργου ισχύος. Ο STATCOM, ως πηγή τάσης, διατηρεί την ικανότητα έγχυσης
      ρεύματος ακόμη και σε χαμηλή τάση δικτύου, ενώ του SVC η απόδοση πέφτει με το
      τετράγωνο της τάσης — γι' αυτό υπερτερεί στη δυναμική στήριξη μετά από σφάλματα.
\item \textbf{Σειριακής σύνδεσης} (ελέγχουν το $X$): ο Ελεγχόμενος με Θυρίστορ Πυκνωτής
      Σειράς (Thyristor Controlled Series Capacitor, TCSC) και ο SSSC. Μεταβάλλοντας τη
      φαινόμενη αντίδραση της γραμμής ανακατευθύνουν τη ροή ισχύος και αυξάνουν τη
      μεταφορική ικανότητα.
\item \textbf{Συνδυασμένης λειτουργίας} (ελέγχουν και τη γωνία): ο Ενοποιημένος Ελεγκτής
      Ροής Ισχύος (Unified Power Flow Controller, UPFC) ελέγχει ταυτόχρονα και ανεξάρτητα
      τάση, ενεργό και άεργο ροή. Είναι ο πληρέστερος και ο ακριβότερος.
\end{itemize}

Σε συνδυασμό με τους συνδέσμους συνεχούς ρεύματος υψηλής τάσης (HVDC), όπου η ροή ισχύος
ορίζεται απευθείας από τον έλεγχο του μετατροπέα, οι διατάξεις FACTS αποτελούν τα βασικά
εργαλεία του Διαχειριστή για τη ρύθμιση τάσης και την ανακούφιση συμφορήσεων στο επίπεδο
μεταφοράς. Τα αντίστοιχα εργαλεία στο επίπεδο διανομής εξετάζονται στα Κεφάλαια 5 και 6.

\subsection{Εκτίμηση κατάστασης}
\label{sec:se}

\subsubsection*{Το πρόβλημα}

Έχουμε δει ότι οι μετρήσεις του SCADA είναι πλεονάζουσες, θορυβώδεις, ασύγχρονες και
ελλιπείς. Έχουμε επίσης δει ότι η ανάλυση ασφάλειας απαιτεί μια \emph{ενιαία και συνεπή}
εικόνα του δικτύου. Η \textbf{εκτίμηση κατάστασης} (state estimation) γεφυρώνει αυτό το
χάσμα. Εισήχθη στα συστήματα ηλεκτρικής ενέργειας από τους Schweppe και Wildes το 1970,
σε μια σειρά τριών εργασιών που παραμένουν μέχρι σήμερα το σημείο αναφοράς, και αποτελεί
έκτοτε τη θεμελιώδη λειτουργία κάθε EMS.

Μια πρώτη σκέψη θα ήταν να λύσουμε απλώς μια ροή ισχύος. Δεν αρκεί, για δύο λόγους. Πρώτον,
η ροή ισχύος απαιτεί \emph{ακριβώς} τα σωστά δεδομένα εισόδου, ενώ εμείς έχουμε πολύ
περισσότερα δεδομένα από όσα χρειάζονται και όλα τους ελαφρώς λανθασμένα. Δεύτερον, η ροή
ισχύος δεν έχει κανέναν τρόπο να ανιχνεύσει ότι κάποιο δεδομένο είναι εσφαλμένο — θα
συγκλίνει ευτυχώς σε μια λάθος λύση. Χρειαζόμαστε μια μέθοδο που \emph{αξιοποιεί} την
πλεονάζουσα πληροφορία αντί να την αγνοεί.

\subsubsection*{Το μοντέλο μέτρησης}

Ως \emph{κατάσταση} ορίζεται το διάνυσμα
$\mathbf{x} = [\boldsymbol{\theta}^{T}, \mathbf{V}^{T}]^{T}$ των γωνιών και των μέτρων των
τάσεων όλων των ζυγών. Για δίκτυο $N$ ζυγών, η γωνία ενός ζυγού ορίζεται αυθαίρετα ως
αναφορά (συνήθως μηδέν), οπότε το πλήθος των αγνώστων είναι

\begin{equation}
n = 2N - 1
\label{eq:nstates}
\end{equation}

Η επιλογή αυτή είναι θεμελιώδης: αν γνωρίζουμε τα $\mathbf{V}$ και $\boldsymbol{\theta}$,
τότε \emph{κάθε} άλλο μέγεθος του δικτύου — ροές, εγχύσεις, απώλειες — υπολογίζεται από τις
εξισώσεις δικτύου. Η κατάσταση είναι το ελάχιστο σύνολο μεγεθών που περιγράφει πλήρως το
σύστημα.

Το μοντέλο μέτρησης είναι

\begin{equation}
\mathbf{z} = \mathbf{h}(\mathbf{x}) + \mathbf{e}
\label{eq:measmodel}
\end{equation}

όπου $\mathbf{z} \in \mathbb{R}^{m}$ το διάνυσμα των μετρήσεων, $\mathbf{h}(\cdot)$ οι μη
γραμμικές συναρτήσεις που συνδέουν την κατάσταση με τα μετρούμενα μεγέθη, και $\mathbf{e}$
ο θόρυβος μέτρησης. Θεωρείται μηδενικής μέσης τιμής, με διαγώνιο πίνακα συνδιακύμανσης
$\mathbf{R} = \mathrm{diag}(\sigma_1^2, \ldots, \sigma_m^2)$ — δηλαδή τα σφάλματα των
διαφορετικών μετρητών θεωρούνται ασυσχέτιστα.

Καθοριστικό μέγεθος είναι ο \textbf{λόγος πλεονασμού}

\begin{equation}
\eta = \frac{m}{n}
\label{eq:redundancy}
\end{equation}

ο οποίος σε πραγματικά συστήματα κυμαίνεται τυπικά μεταξύ 1,7 και 3. Για να πάρουμε μια
αίσθηση μεγεθών: σε δίκτυο $N = 100$ ζυγών έχουμε $n = 199$ αγνώστους, και με $\eta = 2$
περίπου $m = 400$ μετρήσεις. Ο πλεονασμός είναι εκείνος που επιτρέπει τόσο τη μείωση του
θορύβου όσο και — κρίσιμα — την ανίχνευση εσφαλμένων μετρήσεων. Χωρίς αυτόν το πρόβλημα
εκφυλίζεται σε ροή ισχύος, με όλα τα μειονεκτήματα που μόλις αναφέραμε.

\subsubsection*{Σταθμισμένα ελάχιστα τετράγωνα}

Αφού οι μετρήσεις είναι περισσότερες από τους αγνώστους, δεν υπάρχει $\mathbf{x}$ που να τις
ικανοποιεί όλες. Αναζητούμε επομένως εκείνο που τις ικανοποιεί \emph{όσο το δυνατόν
καλύτερα}. Η καθιερωμένη διατύπωση είναι εκείνη των \textbf{σταθμισμένων ελαχίστων
τετραγώνων} (Weighted Least Squares, WLS):

\begin{equation}
J(\mathbf{x}) = \sum_{i=1}^{m} \frac{\left(z_i - h_i(\mathbf{x})\right)^2}{\sigma_i^2}
             = \left[\mathbf{z}-\mathbf{h}(\mathbf{x})\right]^{T} \mathbf{R}^{-1}
               \left[\mathbf{z}-\mathbf{h}(\mathbf{x})\right]
\label{eq:wls}
\end{equation}

Η στάθμιση με $1/\sigma_i^2$ έχει άμεση διαισθητική σημασία: μια μέτρηση με μικρή
τυπική απόκλιση, δηλαδή ακριβής, βαρύνει περισσότερο· μια θορυβώδης βαρύνει λιγότερο.
Ισοδύναμα, κάθε όρος του αθροίσματος είναι το τετράγωνο του σφάλματος \emph{μετρημένο σε
μονάδες της δικής του τυπικής απόκλισης}, ώστε να συγκρίνονται μεταξύ τους μεγέθη
διαφορετικής φύσης και κλίμακας. Υπό την υπόθεση κανονικά κατανεμημένων σφαλμάτων, η
ελαχιστοποίηση της \eqref{eq:wls} ταυτίζεται με την εκτίμηση μέγιστης πιθανοφάνειας.

Επειδή η $\mathbf{h}$ είναι μη γραμμική, η λύση προκύπτει επαναληπτικά με τη μέθοδο
Gauss--Newton: γραμμικοποιούμε γύρω από την τρέχουσα εκτίμηση, λύνουμε το γραμμικό
πρόβλημα, ενημερώνουμε, επαναλαμβάνουμε. Με
$\mathbf{H}(\mathbf{x}) = \partial \mathbf{h}/\partial \mathbf{x}$ τον Ιακωβιανό πίνακα, σε
κάθε επανάληψη $k$ επιλύεται το σύστημα των \emph{κανονικών εξισώσεων}

\begin{align}
\mathbf{G}(\mathbf{x}^{k}) \, \Delta \mathbf{x}^{k}
  &= \mathbf{H}^{T}(\mathbf{x}^{k}) \, \mathbf{R}^{-1}
     \left[\mathbf{z} - \mathbf{h}(\mathbf{x}^{k})\right] \\
\mathbf{G}(\mathbf{x}^{k})
  &= \mathbf{H}^{T}(\mathbf{x}^{k}) \, \mathbf{R}^{-1} \mathbf{H}(\mathbf{x}^{k})
\label{eq:gain}
\end{align}

και ενημερώνεται η εκτίμηση ως $\mathbf{x}^{k+1} = \mathbf{x}^{k} + \Delta \mathbf{x}^{k}$,
έως ότου το $\max|\Delta \mathbf{x}^{k}|$ πέσει κάτω από ένα κατώφλι σύγκλισης. Η σύγκλιση
είναι τυπικά ταχεία, της τάξης των 3--5 επαναλήψεων, επειδή η αρχική εκτίμηση είναι
συνήθως η λύση της προηγούμενης εκτέλεσης — και το σύστημα δεν έχει αλλάξει πολύ στο
μεσοδιάστημα.

Ο πίνακας $\mathbf{G}$ ονομάζεται \emph{πίνακας κέρδους} (gain matrix). Είναι συμμετρικός,
θετικά ημιορισμένος και — κρίσιμα — \emph{αραιός}, επειδή κάθε μέτρηση εμπλέκει μόνο
ελάχιστους ζυγούς. Η αραιότητα δεν είναι λεπτομέρεια υλοποίησης: είναι αυτή που καθιστά
το πρόβλημα επιλύσιμο σε πραγματικό χρόνο για δίκτυα χιλιάδων ζυγών. Μια πυκνή
παραγοντοποίηση θα απαιτούσε χρόνο ανάλογο του $n^3$, ενώ η αραιή κλιμακώνεται σχεδόν
γραμμικά.

\subsubsection*{Παρατηρησιμότητα}

Η επίλυση της \eqref{eq:gain} προϋποθέτει ότι ο $\mathbf{G}$ είναι αντιστρέψιμος. Αυτό δεν
είναι δεδομένο: πολλές μετρήσεις δεν σημαίνει απαραίτητα \emph{κατάλληλες} μετρήσεις. Αν
όλες συγκεντρώνονται σε ένα τμήμα του δικτύου, το υπόλοιπο παραμένει άγνωστο ανεξαρτήτως
πλήθους.

Το σύστημα λέγεται \textbf{παρατηρήσιμο} όταν οι διαθέσιμες μετρήσεις αρκούν για τον
μονοσήμαντο προσδιορισμό ολόκληρης της κατάστασης. Ο έλεγχος γίνεται με δύο ισοδύναμες
προσεγγίσεις:

\begin{itemize}
\item \textbf{Τοπολογική}: αναζητείται δένδρο ζεύξης (spanning tree) του δικτύου πλήρως
      καλυμμένο από μετρήσεις. Είναι γρήγορη και δίνει άμεση φυσική εικόνα.
\item \textbf{Αριθμητική}: ελέγχεται ο βαθμός του $\mathbf{H}$ ή η παραγοντοποίηση του
      $\mathbf{G}$. Είναι πιο γενική και εντοπίζει και οριακές περιπτώσεις.
\end{itemize}

Όταν το σύστημα δεν είναι πλήρως παρατηρήσιμο, εντοπίζονται οι \emph{παρατηρήσιμες νησίδες}
και η κατάσταση εκτιμάται μόνο εντός αυτών. Τα κενά γεφυρώνονται με \emph{ψευδομετρήσεις}
(pseudo-measurements) — εκτιμήσεις φορτίου από ιστορικά δεδομένα ή προβλέψεις, στις οποίες
αποδίδεται μεγάλη τυπική απόκλιση ώστε να βαρύνουν ελάχιστα.

\subsubsection*{Ανίχνευση και αναγνώριση εσφαλμένων μετρήσεων}

Εδώ αποδίδει ο πλεονασμός. Αν μια μέτρηση είναι εσφαλμένη, θα «διαφωνεί» με τις υπόλοιπες,
και η διαφωνία αυτή αποτυπώνεται στα \emph{υπόλοιπα}
$\mathbf{r} = \mathbf{z}-\mathbf{h}(\hat{\mathbf{x}})$. Η κλασική διαδικασία έχει δύο στάδια:

\textbf{Στάδιο 1 — Ανίχνευση.} Υπό την υπόθεση ότι όλες οι μετρήσεις είναι σωστές, η
αντικειμενική συνάρτηση $J(\hat{\mathbf{x}})$ στη λύση ακολουθεί κατανομή $\chi^2$ με
$m-n$ βαθμούς ελευθερίας. Στο παράδειγμά μας με $m = 400$ και $n = 199$, οι βαθμοί
ελευθερίας είναι 201 και η αναμενόμενη τιμή του $J$ είναι επίσης περίπου 201. Αν η
υπολογισμένη τιμή υπερβεί το κατώφλι $\chi^2_{m-n,\,\alpha}$ για δεδομένο επίπεδο
εμπιστοσύνης $\alpha$, τότε υπάρχει στατιστική ένδειξη ότι κάτι δεν πάει καλά. Το τεστ
όμως λέει μόνο \emph{ότι} υπάρχει πρόβλημα, όχι \emph{πού}.

\textbf{Στάδιο 2 — Αναγνώριση.} Τα ακατέργαστα υπόλοιπα δεν είναι συγκρίσιμα μεταξύ τους,
γιατί κάθε μέτρηση έχει διαφορετική ακρίβεια και διαφορετικό βαθμό «στήριξης» από τις
γειτονικές. Κανονικοποιούνται επομένως ως

\begin{equation}
r_i^{N} = \frac{|r_i|}{\sqrt{\Omega_{ii}}}, \qquad
\boldsymbol{\Omega} = \mathbf{R} - \mathbf{H}\mathbf{G}^{-1}\mathbf{H}^{T}
\label{eq:normres}
\end{equation}

όπου $\boldsymbol{\Omega}$ ο πίνακας συνδιακύμανσης των υπολοίπων. Η μέτρηση με το μέγιστο
$r_i^{N}$ — κριτήριο Largest Normalized Residual — θεωρείται η ύποπτη, απορρίπτεται, και η
εκτίμηση επαναλαμβάνεται χωρίς αυτήν. Η διαδικασία επαναλαμβάνεται έως ότου το τεστ $\chi^2$
περάσει. Ένα κατώφλι $r_i^{N} > 3$ αντιστοιχεί σε γεγονός τριών τυπικών αποκλίσεων και
χρησιμοποιείται συχνά στην πράξη.

\subsubsection*{Δύο θεμελιώδεις περιορισμοί}

Η μέθοδος είναι ισχυρή, αλλά έχει δύο όρια που πρέπει να είναι ξεκάθαρα.

\textbf{Πρώτον, οι κρίσιμες μετρήσεις.} Μια μέτρηση λέγεται \emph{κρίσιμη} όταν η αφαίρεσή
της καθιστά το σύστημα μη παρατηρήσιμο. Τέτοιες μετρήσεις έχουν εξ ορισμού μηδενικό
υπόλοιπο — δεν υπάρχει άλλη πληροφορία για να τις διαψεύσει, οπότε η εκτίμηση αναγκάζεται
να τις αναπαραγάγει ακριβώς. Συνεπώς ένα σφάλμα σε κρίσιμη μέτρηση είναι \emph{μη
ανιχνεύσιμο}, όσο μεγάλο κι αν είναι. Ο εντοπισμός των κρίσιμων μετρήσεων ενός δικτύου
είναι επομένως μέρος του σχεδιασμού του συστήματος μέτρησης, όχι ακαδημαϊκή άσκηση.

\textbf{Δεύτερον, οι επιθέσεις έγχυσης ψευδών δεδομένων.} Οι Liu, Ning και Reiter έδειξαν
το 2009 κάτι πιο ανησυχητικό. Έστω επιτιθέμενος που γνωρίζει την τοπολογία και τις
παραμέτρους του δικτύου, δηλαδή τον πίνακα $\mathbf{H}$. Αν αλλοιώσει τις μετρήσεις κατά
$\mathbf{a} = \mathbf{H}\mathbf{c}$ για οποιοδήποτε διάνυσμα $\mathbf{c}$, τότε η
αλλοιωμένη μέτρηση $\mathbf{z}_{\text{bad}} = \mathbf{z} + \mathbf{a}$ αντιστοιχεί ακριβώς
στην κατάσταση $\hat{\mathbf{x}} + \mathbf{c}$, και το υπόλοιπο παραμένει \emph{αμετάβλητο}:

\begin{equation}
\mathbf{r}_{\text{bad}} = \mathbf{z} + \mathbf{H}\mathbf{c}
        - \mathbf{H}(\hat{\mathbf{x}} + \mathbf{c}) = \mathbf{r}
\label{eq:fdia}
\end{equation}

Η επίθεση διαφεύγει πλήρως των παραπάνω στατιστικών ελέγχων, επειδή είναι \emph{συνεπής} με
το μοντέλο. Δεν πρόκειται για ελάττωμα υλοποίησης αλλά για δομικό όριο κάθε ελέγχου που
βασίζεται στα υπόλοιπα. Οι επιθέσεις αυτές, γνωστές ως False Data Injection Attacks (FDIA),
και οι στρατηγικές αντιμετώπισής τους εξετάζονται αναλυτικά στο Κεφάλαιο 8.

\subsubsection*{Ενσωμάτωση μετρήσεων PMU}

Η προσθήκη μετρήσεων PMU μεταβάλλει ουσιωδώς το πρόβλημα. Καθώς η PMU μετρά απευθείας τη
γωνία τάσης, η σχέση μεταξύ μετρήσεων και κατάστασης γίνεται \emph{γραμμική} όταν εκφραστεί
σε ορθογώνιες συντεταγμένες. Αν το σύστημα είναι πλήρως παρατηρήσιμο μόνο από PMU, η
\eqref{eq:measmodel} εκφυλίζεται σε $\mathbf{z} = \mathbf{H}\mathbf{x} + \mathbf{e}$ και η
λύση δίνεται σε ένα βήμα:

\begin{equation}
\hat{\mathbf{x}} = \left(\mathbf{H}^{T}\mathbf{R}^{-1}\mathbf{H}\right)^{-1}
                    \mathbf{H}^{T}\mathbf{R}^{-1}\mathbf{z}
\label{eq:linse}
\end{equation}

Τα πλεονεκτήματα είναι σημαντικά: δεν υπάρχουν επαναλήψεις, άρα ούτε προβλήματα σύγκλισης
ή εξάρτηση από την αρχική εκτίμηση· ο $\mathbf{H}$ είναι σταθερός για δεδομένη τοπολογία,
άρα προϋπολογίζεται· και η εκτίμηση μπορεί να εκτελείται με τον ρυθμό των PMU, δηλαδή
δεκάδες φορές ανά δευτερόλεπτο αντί για μία φορά ανά λεπτό.

Στην πράξη, η πλήρης κάλυψη με PMU σπανίζει για λόγους κόστους, οπότε χρησιμοποιούνται
\textbf{υβριδικοί εκτιμητές} που συνδυάζουν συμβατικές μετρήσεις SCADA με συγχρονοφασιθέτες.
Η βασική δυσκολία εκεί είναι η ετερογένεια των ρυθμών: πώς συνδυάζονται μετρήσεις που
ανανεώνονται κάθε 4 s με άλλες που ανανεώνονται κάθε 20 ms.

\subsection{Στάδια λειτουργίας: από τον προγραμματισμό στον πραγματικό χρόνο}

Η λειτουργία του συστήματος δεν αποφασίζεται σε μία στιγμή, αλλά μέσα από διαδοχικές
χρονικές κλίμακες. Η λογική είναι απλή και αξίζει να διατυπωθεί ρητά: \emph{όσο πλησιάζει η
ώρα παράδοσης, τόσο ακριβέστερη γίνεται η πρόβλεψη — και τόσο ακριβότερη η διόρθωση}. Κάθε
στάδιο διορθώνει τα σφάλματα του προηγούμενου, με λιγότερο χρόνο και λιγότερες διαθέσιμες
επιλογές. Στο ευρωπαϊκό μοντέλο αγοράς (Target Model) οι κλίμακες αυτές αντιστοιχούν σε
διακριτές αγορές.

\begin{itemize}
\item \textbf{Μακροπρόθεσμος και προθεσμιακός ορίζοντας} (έτη έως εβδομάδες). Μελέτες
      επάρκειας ισχύος, συντονισμός προγραμμάτων συντήρησης γραμμών και μονάδων, εκτίμηση
      των αναγκών σε εφεδρείες, προθεσμιακά προϊόντα αντιστάθμισης κινδύνου.
\item \textbf{Αγορά Επόμενης Ημέρας} (day-ahead). Το κύριο στάδιο προγραμματισμού. Οι
      συμμετέχοντες υποβάλλουν προσφορές για κάθε αγοραία περίοδο της επόμενης ημέρας και η
      αγορά εκκαθαρίζεται με ενιαία σύζευξη σε ευρωπαϊκό επίπεδο (Single Day-Ahead Coupling),
      ώστε η ενέργεια να ρέει από τις φθηνότερες προς τις ακριβότερες περιοχές μέχρι να
      κορεστεί η διαθέσιμη διασυνοριακή ικανότητα. Ο Διαχειριστής ελέγχει στη συνέχεια αν
      το προκύπτον πρόγραμμα είναι τεχνικά εφικτό.
\item \textbf{Ενδοημερήσια Αγορά} (intra-day). Επιτρέπει τη διόρθωση των θέσεων καθώς
      πλησιάζει η ώρα παράδοσης και βελτιώνονται οι προβλέψεις. Έχει ιδιαίτερη σημασία σε
      συστήματα με υψηλή διείσδυση αιολικής και φωτοβολταϊκής παραγωγής: το σφάλμα πρόβλεψης
      αιολικής παραγωγής μειώνεται δραστικά όσο μικραίνει ο ορίζοντας, οπότε η δυνατότητα
      αναπροσαρμογής λίγες ώρες πριν την παράδοση έχει μεγάλη οικονομική αξία.
\item \textbf{Πραγματικός χρόνος και εξισορρόπηση} (balancing). Ο Διαχειριστής ενεργοποιεί
      προσφορές ενέργειας εξισορρόπησης ώστε να καλύψει τις αποκλίσεις που απομένουν, να
      αντιμετωπίσει διαταραχές και να διαχειριστεί συμφορήσεις.
\end{itemize}

\subsubsection*{Η αγοραία περίοδος}

Η \emph{αγοραία περίοδος} (Market Time Unit, MTU) είναι η χρονική ανάλυση με την οποία
εκκαθαρίζεται η αγορά. Για χρόνια ήταν η ώρα. Από τις 30 Σεπτεμβρίου 2025 (ημέρα παράδοσης
1η Οκτωβρίου 2025) η ενιαία σύζευξη επόμενης ημέρας λειτουργεί με MTU των \textbf{15 λεπτών}
σε όλες τις ευρωπαϊκές ζώνες προσφοράς, εναρμονισμένη με την περίοδο εκκαθάρισης
αποκλίσεων. Η αλλαγή δεν είναι τεχνική λεπτομέρεια: μικρότερη αγοραία περίοδος σημαίνει ότι
οι απότομες μεταβολές παραγωγής και φορτίου αποτυπώνονται στο ίδιο το πρόγραμμα αντί να
μετακυλίονται στην εξισορρόπηση, μειώνοντας τις ανάγκες σε εφεδρείες.

\subsubsection*{Το ελληνικό πλαίσιο}

Στην Ελλάδα το μοντέλο τέθηκε σε ισχύ την 1η Νοεμβρίου 2020. Η Αγορά Επόμενης Ημέρας και η
Ενδοημερήσια Αγορά λειτουργούν από το Ελληνικό Χρηματιστήριο Ενέργειας (ΕΧΕ), ενώ την
\textbf{Αγορά Εξισορρόπησης} λειτουργεί ο ΑΔΜΗΕ. Ο τελευταίος εκτελεί τη Διαδικασία
Ολοκληρωμένου Προγραμματισμού και ενεργοποιεί σε πραγματικό χρόνο τόσο ενέργεια
εξισορρόπησης όσο και τις εφεδρείες που περιγράφονται στη συνέχεια.

\subsection{Απόκριση συχνότητας και εφεδρείες}

\subsubsection*{Αδράνεια και ο αρχικός ρυθμός}

Η συχνότητα είναι ο πλέον άμεσος δείκτης του ισοζυγίου ισχύος: παραμένει σταθερή όσο
παραγωγή και ζήτηση εξισορροπούνται και αποκλίνει μόλις εμφανιστεί ανισοζύγιο.

Αμέσως μετά από μια διαταραχή ισχύος $\Delta P$ — τυπικά την απώλεια μιας μεγάλης μονάδας —
κανένας ρυθμιστής δεν έχει προλάβει να αντιδράσει. Το έλλειμμα καλύπτεται από την κινητική
ενέργεια των στρεφόμενων μαζών, οι οποίες επιβραδύνονται. Ο αρχικός ρυθμός μεταβολής της
συχνότητας (Rate of Change of Frequency, ROCOF) προκύπτει από την εξίσωση ταλάντωσης και
είναι

\begin{equation}
\frac{\mathrm{d}f}{\mathrm{d}t} = \frac{f_0 \, \Delta P}{2 H S_n}
\label{eq:rocof}
\end{equation}

όπου $f_0$ η ονομαστική συχνότητα, $H$ η σταθερά αδράνειας του συστήματος σε δευτερόλεπτα
και $S_n$ η συνολική ονομαστική φαινόμενη ισχύς των \emph{συγχρονισμένων} μονάδων.

\emph{Παράδειγμα.} Για σύγχρονη περιοχή με $S_n = 300$ GVA, $H = 5$ s και απώλεια
$\Delta P = 3\,000$ MW, ο αρχικός ρυθμός είναι
$50 \times 3000 / (2 \times 5 \times 300\,000) = 0{,}05$ Hz/s. Αν η ίδια απώλεια συμβεί σε
ώρα χαμηλού φορτίου με τη μισή συγχρονισμένη ισχύ, ο ρυθμός διπλασιάζεται.

Η σχέση \eqref{eq:rocof} εξηγεί γιατί η αντικατάσταση σύγχρονων γεννητριών από μονάδες
συνδεδεμένες μέσω μετατροπέων — οι οποίες δεν προσφέρουν εγγενή αδράνεια — οδηγεί σε
ταχύτερες αποκλίσεις συχνότητας. Ο χρόνος που μεσολαβεί μέχρι η συχνότητα να φτάσει σε
επικίνδυνα επίπεδα συρρικνώνεται, και η απόκριση πρέπει να γίνει αντίστοιχα ταχύτερη. Από
εδώ πηγάζει το ενδιαφέρον για συνθετική αδράνεια και ταχεία απόκριση συχνότητας.

\subsubsection*{Στατισμός και πρωτεύουσα ρύθμιση}

Η πρωτεύουσα ρύθμιση υλοποιείται από τους ρυθμιστές στροφών των μονάδων, οι οποίοι
μεταβάλλουν την παραγωγή ανάλογα προς την απόκλιση συχνότητας. Ο συντελεστής αναλογίας
εκφράζεται μέσω του \emph{στατισμού} $R$:

\begin{equation}
\frac{\Delta P}{P_n} = -\frac{1}{R}\,\frac{\Delta f}{f_0}
\label{eq:droop}
\end{equation}

Τυπικές τιμές είναι $R = 4$--5\%. \emph{Παράδειγμα:} μονάδα 500 MW με $R = 5\%$, σε
απόκλιση $\Delta f = -200$ mHz (δηλαδή $-0{,}4\%$), αυξάνει την παραγωγή της κατά
$(0{,}004/0{,}05) \times 500 = 40$ MW.

Το αρνητικό πρόσημο εξασφαλίζει αρνητική ανάδραση: πτώση συχνότητας προκαλεί αύξηση
παραγωγής. Το ουσιώδες όμως είναι ότι η ρύθμιση είναι \emph{αναλογική} και όχι
ολοκληρωτική. Για να παραμείνει η επιπλέον παραγωγή, πρέπει να παραμείνει και η απόκλιση
συχνότητας. Η πρωτεύουσα ρύθμιση επομένως \textbf{σταθεροποιεί} τη συχνότητα σε μια νέα
μόνιμη τιμή, δεν την \textbf{επαναφέρει} στα 50 Hz. Ακριβώς αυτό το υπόλοιπο σφάλμα καλείται
να εξαλείψει το επόμενο επίπεδο.

\subsubsection*{Δευτερεύουσα ρύθμιση και το σφάλμα ελέγχου περιοχής}

Η δευτερεύουσα ρύθμιση εκτελείται από τον κεντρικό ελεγκτή κάθε Διαχειριστή, ο οποίος
επιδιώκει τον μηδενισμό του \emph{Σφάλματος Ελέγχου Περιοχής} (Area Control Error):

\begin{equation}
\mathrm{ACE} = \Delta P_{\text{ανταλλαγών}} + K \, \Delta f
\label{eq:ace}
\end{equation}

όπου $\Delta P_{\text{ανταλλαγών}}$ η απόκλιση των διασυνοριακών ανταλλαγών από το πρόγραμμα
και $K$ ο συντελεστής της περιοχής. Ο συνδυασμός των δύο όρων είναι εύγλωττος: αν η
διαταραχή συνέβη \emph{εντός} της περιοχής, και οι δύο όροι είναι μη μηδενικοί και το ACE
είναι μεγάλο, οπότε η περιοχή διορθώνει. Αν συνέβη \emph{εκτός}, οι δύο όροι τείνουν να
αλληλοαναιρούνται και το ACE παραμένει μικρό — η περιοχή έχει ήδη προσφέρει πρωτεύουσα
ρύθμιση και δεν χρειάζεται να κάνει περισσότερα. Έτσι κάθε Διαχειριστής διορθώνει τα δικά
του ανισοζύγια, ενώ η αλληλεγγύη στην πρώτη στιγμή παραμένει καθολική.

\subsubsection*{Τα τέσσερα επίπεδα}

\begin{itemize}
\item \textbf{FCR} --- Εφεδρεία Διατήρησης Συχνότητας (πρωτεύουσα ρύθμιση). Αυτόματη και
      τοπική, ενεργοποιείται από όλες τις μονάδες της σύγχρονης περιοχής από κοινού.
      Σταθεροποιεί τη συχνότητα.
\item \textbf{aFRR} --- Αυτόματη Εφεδρεία Αποκατάστασης Συχνότητας (δευτερεύουσα ρύθμιση).
      Αυτόματη και κεντρική, οδηγείται από το ACE. Επαναφέρει τη συχνότητα στα 50 Hz και
      απελευθερώνει την FCR.
\item \textbf{mFRR} --- Χειροκίνητη Εφεδρεία Αποκατάστασης Συχνότητας (τριτεύουσα ρύθμιση).
      Ενεργοποιείται με εντολή του Διαχειριστή και αντικαθιστά την aFRR σε παρατεταμένα
      ανισοζύγια.
\item \textbf{RR} --- Εφεδρεία Αντικατάστασης. Αποκαθιστά τα επίπεδα των προηγούμενων
      εφεδρειών ώστε το σύστημα να είναι έτοιμο για επόμενη διαταραχή.
\end{itemize}

Παρατηρήστε το μοτίβο: κάθε επίπεδο είναι βραδύτερο, φθηνότερο και πιο στοχευμένο από το
προηγούμενο, και ο ρόλος του είναι να \emph{απελευθερώσει} το ταχύτερο ώστε αυτό να είναι
και πάλι διαθέσιμο.

\begin{table}[!htbp]
\centering
\begin{tabular}{lllc}
\hline
Εφεδρεία & Ονομασία & Ενεργοποίηση & Πλήρης ενεργοποίηση \\
\hline
FCR  & Πρωτεύουσα   & αυτόματη, τοπική  & 30 s \\
aFRR & Δευτερεύουσα & αυτόματη, κεντρική & 5 min \\
mFRR & Τριτεύουσα   & χειροκίνητη        & 12,5 min \\
RR   & Αντικατάστασης & χειροκίνητη      & $>$ 15 min \\
\hline
\end{tabular}
\caption{Επίπεδα εφεδρείας συχνότητας κατά τον SO GL. Οι χρόνοι πλήρους ενεργοποίησης
αφορούν τη σύγχρονη περιοχή της Ηπειρωτικής Ευρώπης.}
\label{tab:reserves}
\end{table}

\subsubsection*{Ποσοτικές απαιτήσεις για την Ηπειρωτική Ευρώπη}

\begin{itemize}
\item Το \textbf{συμβάν αναφοράς} για τη διαστασιολόγηση της FCR είναι απώλεια
      3\,000 MW, και αντίστοιχα 3\,000 MW στην αντίθετη κατεύθυνση. Η ελάχιστη διαθέσιμη FCR
      στη σύγχρονη περιοχή είναι ανά πάσα στιγμή 3\,000 MW.
\item Η μέγιστη απόκλιση συχνότητας μόνιμης κατάστασης για το συμβάν αναφοράς είναι
      \textbf{200 mHz}. Από τα δύο αυτά μεγέθη προκύπτει ο συντελεστής αναλογίας μεταξύ FCR
      και απόκλισης συχνότητας: $3\,000\ \text{MW} / 0{,}2\ \text{Hz} = 15\,000$ MW/Hz.
\item Το \textbf{τυπικό εύρος συχνότητας} είναι $\pm 50$ mHz και δεν επιτρέπεται να
      παραβιάζεται για περισσότερα από 15\,000 λεπτά ανά έτος. Η μέγιστη στιγμιαία απόκλιση
      ορίζεται σε 800 mHz.
\item Απόκλιση συχνότητας μεγαλύτερη των 200 mHz οδηγεί σε κατάσταση \emph{έκτακτης ανάγκης}.
\end{itemize}

\subsection{Σύνοψη}

Η λειτουργία του ΣΜΗΕ στηρίζεται στον κλειστό βρόχο με τον οποίο ξεκινήσαμε:

\begin{enumerate}
\item Η \textbf{υποδομή εποπτείας} (SCADA, PMU) παράγει μετρήσεις — πλεονάζουσες,
      θορυβώδεις και εν μέρει ασύγχρονες.
\item Η \textbf{εκτίμηση κατάστασης} τις μετατρέπει σε μία συνεπή και στατιστικά βέλτιστη
      εικόνα του δικτύου, απορρίπτοντας ταυτόχρονα τα προφανώς εσφαλμένα δεδομένα.
\item Τα \textbf{εργαλεία ανάλυσης ασφάλειας} αξιολογούν την εικόνα αυτή έναντι του
      κριτηρίου N-1 και κατατάσσουν το σύστημα σε μία από τις πέντε καταστάσεις λειτουργίας.
\item Οι \textbf{αποφάσεις} υλοποιούνται μέσω των αγορών και των διατάξεων ελέγχου, σε
      κλίμακες που εκτείνονται από τα δευτερόλεπτα της πρωτεύουσας ρύθμισης έως τους μήνες
      του προγραμματισμού.
\end{enumerate}

Ο βρόχος είναι τόσο ισχυρός όσο ο ασθενέστερος κρίκος του — και είδαμε ότι ο κρίκος αυτός
είναι συχνά η αξιοπιστία της ίδιας της εικόνας: κρίσιμες μετρήσεις χωρίς πλεονασμό,
επιθέσεις συνεπείς με το μοντέλο, τμήματα του δικτύου που δεν είναι παρατηρήσιμα.

Οι υπηρεσίες που στηρίζουν αυτόν τον βρόχο ονομάζονται \emph{επικουρικές υπηρεσίες} και
διακρίνονται σε σχετικές και μη σχετικές με τη συχνότητα· τις εξετάζουν αντίστοιχα τα
Κεφάλαια 3 και 2. Το αντίστοιχο πρόβλημα στο επίπεδο της διανομής, όπου η παρατηρησιμότητα
είναι πολύ χαμηλότερη και οι ψευδομετρήσεις κυριαρχούν, εξετάζεται στο Κεφάλαιο 4. Η
ασφάλεια της ίδιας της υποδομής μέτρησης και ελέγχου αποτελεί το αντικείμενο του
Κεφαλαίου 8.

\section{Παραδείγματα}
\label{kef-01-paradeigmata}
% TODO: παραδείγματα και λυμένες ασκήσεις

\section{Βιβλιογραφία}
\label{kef-01-vivliografia}

\subsubsection*{Βιβλία αναφοράς}

\begin{itemize}
\item A. Abur and A. Gómez Expósito, \textit{Power System State Estimation: Theory and
      Implementation}. New York: Marcel Dekker, 2004.
\item N. G. Hingorani and L. Gyugyi, \textit{Understanding FACTS: Concepts and Technology of
      Flexible AC Transmission Systems}. New York: IEEE Press, 2000.
\item P. Kundur, \textit{Power System Stability and Control}. New York: McGraw-Hill, 1994.
\item A. Monticelli, \textit{State Estimation in Electric Power Systems: A Generalized
      Approach}. Boston, MA: Kluwer Academic Publishers, 1999.
\item A. G. Phadke and J. S. Thorp, \textit{Synchronized Phasor Measurements and Their
      Applications}, 2nd ed. Cham: Springer, 2017.
\item A. J. Wood, B. F. Wollenberg and G. B. Sheblé, \textit{Power Generation, Operation, and
      Control}, 3rd ed. Hoboken, NJ: Wiley, 2014.
\end{itemize}

\subsubsection*{Θεμελιώδεις εργασίες}

\begin{itemize}
\item T. E. Dy Liacco, «The Adaptive Reliability Control System», \textit{IEEE Transactions
      on Power Apparatus and Systems}, vol. PAS-86, no. 5, pp. 517--531, May 1967.
      doi: \texttt{10.1109/TPAS.1967.291728}
\item Y. Liu, P. Ning and M. K. Reiter, «False Data Injection Attacks against State
      Estimation in Electric Power Grids», \textit{ACM Transactions on Information and System
      Security}, vol. 14, no. 1, art. 13, May 2011.
      doi: \texttt{10.1145/1952982.1952995}
\item A. Monticelli, «Electric Power System State Estimation», \textit{Proceedings of the
      IEEE}, vol. 88, no. 2, pp. 262--282, Feb. 2000.
      doi: \texttt{10.1109/5.824004}
\item F. C. Schweppe and J. Wildes, «Power System Static-State Estimation, Part I: Exact
      Model», \textit{IEEE Transactions on Power Apparatus and Systems}, vol. PAS-89, no. 1,
      pp. 120--125, Jan. 1970. doi: \texttt{10.1109/TPAS.1970.292678}
\item F. C. Schweppe and D. B. Rom, «Power System Static-State Estimation, Part II:
      Approximate Model», \textit{IEEE Transactions on Power Apparatus and Systems},
      vol. PAS-89, no. 1, pp. 125--130, Jan. 1970.
\item F. C. Schweppe, «Power System Static-State Estimation, Part III: Implementation»,
      \textit{IEEE Transactions on Power Apparatus and Systems}, vol. PAS-89, no. 1,
      pp. 130--135, Jan. 1970.
\end{itemize}

\subsubsection*{Πρότυπα}

\begin{itemize}
\item IEC/IEEE 60255-118-1:2018, \textit{Measuring Relays and Protection Equipment ---
      Part 118-1: Synchrophasor for Power Systems --- Measurements}. IEC/IEEE, 2018.
\item IEC 61850, \textit{Communication Networks and Systems for Power Utility Automation}. IEC.
\item IEC 60870-5-104, \textit{Telecontrol Equipment and Systems --- Part 5-104: Network
      Access for IEC 60870-5-101 Using Standard Transport Profiles}. IEC.
\end{itemize}

\subsubsection*{Κανονιστικό πλαίσιο και Διαχειριστές}

\begin{itemize}
\item Commission Regulation (EU) 2017/1485 of 2 August 2017 establishing a guideline on
      electricity transmission system operation (System Operation Guideline, SO GL).
      \textit{Official Journal of the European Union}, L 220, 25.8.2017.
      \href{https://eur-lex.europa.eu/eli/reg/2017/1485/oj}{eur-lex.europa.eu/eli/reg/2017/1485/oj}
\item Commission Regulation (EU) 2017/2195 of 23 November 2017 establishing a guideline on
      electricity balancing (Electricity Balancing Guideline, EB GL). \textit{Official Journal
      of the European Union}, L 312, 28.11.2017.
      \href{https://eur-lex.europa.eu/eli/reg/2017/2195/oj}{eur-lex.europa.eu/eli/reg/2017/2195/oj}
\item ENTSO-E, \textit{Synchronous Area Framework Agreement for Regional Group Continental
      Europe --- Policy 1: Load-Frequency Control and Reserves}, 2019.
      \href{https://www.entsoe.eu/}{entsoe.eu}
\item ΑΔΜΗΕ Α.Ε., \textit{Κώδικας Διαχείρισης του Ελληνικού Συστήματος Μεταφοράς Ηλεκτρικής
      Ενέργειας}. \href{https://www.admie.gr/}{admie.gr}
\item ΑΔΜΗΕ Α.Ε., \textit{Κανονισμός Αγοράς Εξισορρόπησης}. \href{https://www.admie.gr/}{admie.gr}
\item Ελληνικό Χρηματιστήριο Ενέργειας (ΕΧΕ), \textit{Κανονισμός Αγοράς Επόμενης Ημέρας και
      Ενδοημερήσιας Αγοράς}. \href{https://www.enexgroup.gr/}{enexgroup.gr}
\end{itemize}
